\section{总结与延伸}

\subsection{核心要点}

\begin{enumerate}
    \item 主动交互是智能的基础——纯文本 LLM 是``被动小猫''
    \item Foundation Agent（Voyager、Eureka）展示了 LLM 驱动的开放世界探索和自动 reward 设计
    \item Project GR00T 旨在构建跨机器人形态的通用基础模型
    \item 数据稀缺和 sim-to-real gap 是核心挑战
    \item 具身 AI 是 AGI 的必要组成部分
\end{enumerate}

\subsection{Summary table}

\begin{table}[H]
\centering
\begin{tabular}{p{0.24\textwidth} p{0.2\textwidth} p{0.28\textwidth} p{0.2\textwidth}}
\toprule
主题 & 承诺 & 挑战 & 关键产出 \\
\midrule
主动具身交互 & 定制机器人感知-语言-动作闭环 & 数据稀缺 + sim-to-real gap & Voyager/Eureka 架构、GR00T baseline \\
数据与训练 & LLM-generated tasks + replay hygiene & 多样性保持、自然分布 & Isaac Sim pipelines、LLM reward generation \\
部署运营 & Trace + metrics + alerts & 工具更换、calibration drift & Incident playback、checklists \\
治理 & Safety board + audit trail & Autonomous failure modes & Usage policy + override switch \\
\bottomrule
\end{tabular}
\caption{Lecture 04 的主要落地纬度}
\end{table}

\subsection{拓展阅读}

\begin{itemize}
    \item Wang et al., ``Voyager: An Open-Ended Embodied Agent with Large Language Models,'' 2023.
    \item Ma et al., ``Eureka: Human-Level Reward Design via Coding Large Language Models,'' ICLR 2024.
    \item Fan et al., ``MineDojo: Building Open-Ended Embodied Agents with Internet-Scale Knowledge,'' NeurIPS 2022.
\end{itemize}


\end{document}
